
\documentclass{article}
%%%%%%%%%%%%%%%%%%%%%%%%%%%%%%%%%%%%%%%%%%%%%%%%%%%%%%%%%%%%%%%%%%%%%%%%%%%%%%%%%%%%%%%%%%%%%%%%%%%%%%%%%%%%%%%%%%%%%%%%%%%%%%%%%%%%%%%%%%%%%%%%%%%%%%%%%%%%%%%%%%%%%%%%%%%%%%%%%%%%%%%%%%%%%%%%%%%%%%%%%%%%%%%%%%%%%%%%%%%%%%%%%%%%%%%%%%%%%%%%%%%%%%%%%%%%
\usepackage{amsmath}
\usepackage{chicago}
\usepackage{setspace}
\usepackage{geometry}
\usepackage{graphicx}

\setcounter{MaxMatrixCols}{10}

\newtheorem{theorem}{Theorem}
\newtheorem{acknowledgement}[theorem]{Acknowledgement}
\newtheorem{algorithm}[theorem]{Algorithm}
\newtheorem{axiom}[theorem]{Axiom}
\newtheorem{case}[theorem]{Case}
\newtheorem{claim}[theorem]{Claim}
\newtheorem{conclusion}[theorem]{Conclusion}
\newtheorem{condition}[theorem]{Condition}
\newtheorem{conjecture}[theorem]{Conjecture}
\newtheorem{corollary}[theorem]{Corollary}
\newtheorem{criterion}[theorem]{Criterion}
\newtheorem{definition}[theorem]{Definition}
\newtheorem{example}[theorem]{Example}
\newtheorem{exercise}[theorem]{Exercise}
\newtheorem{lemma}[theorem]{Lemma}
\newtheorem{notation}[theorem]{Notation}
\newtheorem{problem}[theorem]{Problem}
\newtheorem{proposition}[theorem]{Proposition}
\newtheorem{remark}[theorem]{Remark}
\newtheorem{solution}[theorem]{Solution}
\newtheorem{summary}[theorem]{Summary}
\newenvironment{proof}[1][Proof]{\noindent\textbf{#1.} }{\ \rule{0.5em}{0.5em}}
\geometry{margin=1in}

\begin{document}

\title{The Many Faces of the Taylor Rule for Advanced Undergraduate
Macroeconomics}
\date{\textsf{May 5, 2022}}
\author{George A. Waters \\
%EndAName
Department of Economics\\
Illinois State University\\
Normal, IL 61790-4200}
\maketitle

\begin{abstract}
The Taylor Rule and Fisher Relation can be represented on a graph that
allows for discussion of the zero lower bound on interest rates, the
existence of multiple equilibria, secular stagnation and Japan's lost
decade, among other issues. The Taylor Rule and Fisher Relation can also be
included in a small macroeconomic model that can be used to study the
stability under various interest rate rules under adaptive expectations or
solved under rational expectations. \ There are a number of related
empirical exercises using both descriptive statistics and regressions for
undergraduates with intermediate macro level knowledge.

\textbf{Keywords: \ }Monetary Policy, Interest Rate Rules, Taylor Rule,
zero-lower bound

\textbf{JEL classification: }A22, E0, E52, E60
\end{abstract}

\begin{doublespace}
\pagebreak

\section{Introduction}

Taylor Rules are a ubiquitous method for describing the current conduct of
monetary policy. \ At the same time, the zero lower bound on interest rates
(ZLB) is a key constraint on policy and has implications the type of
stabilization policies that work. \ A combination of the Fisher relation and
an interest rate rule that accounts for the ZLB on a graph or as part of an
algebraic model offers insights on these issues that are accessible to
undergraduates. In addition, there are a number of theoretical issues and
empirical exercises that arise from extensions.

An important point that can be made using a model with a Taylor Rule (TR)
and the Fisher relation (FR) concern the existence of multiple equilibria
and their stability. The equilibrium at the inflation target is stable under
learning if the Taylor Principle is satisfied; however, equilibria at an
interest rate peg, such as the ZLB, are not, under reasonable assumptions.
Furthermore, a graph with the TR and FR provides intuition about the impact
of secular stagnation, changes in the inflation target and Neo-Fisherian
increases in the lower bound on interest rates.

Including standard loanable funds and IS-LM graphs in the analysis allows
for a discussion of the excess savings under the ZLB in financial markets
and the liquidity trap in the broader economy. The liquidity trap is the
primary counterexample to Say's Law, and Keynes' (1935) argument for the
necessity of fiscal stimulus in such a situation is apparent. The
extraordinary policies of forward guidance and quantitative easing employed
by the Federal Reserve in response to the Great Recession of 2009-10 can be
interpreted on the IS-LM graph as well.

An algebraic model with IS and aggregate supply relations included with the
TR and FR can be analyzed under different assumptions about expectations.
Solving the model under rational expectations shows the role of the
parameters in the TR on the response to shocks. Under adaptive expectations,
the dynamics of the resulting difference equation(s) demonstrates the
importance of the Taylor Principle and the instability arising from interest
rate pegs\footnote{%
This model is closely related to the graphical analysis in Buttet and Roy
(2014).}.

Another advantage to focusing on the TR and FR is the close relation to the
data, exemplified in the graphs of the recent experience of the U.S. and
Japan in Bullard (2010). There are a number of empirical exercises students
can do such as comparing graphs of the recommended interest rates for
various Taylor rules and the actual policy rates. Estimating interest rate
rules using the methods of Judd and Rudebusch (1998) is an excellent way to
get intuition about policy and experience applying econometric techniques.

Finally, analyzing the graphical and algebraic models described here forms
and excellent basis for students planning to study these issues at a
graduate level. The TR-FR graphs arise often in theoretical papers, see
footnote 3 for examples. \ Solving under rational expectations with the
method of undetermined coefficients and analyzing stability under learning
are methods that are invaluable for further study.

The conduct of monetary policy at the Federal Reserve is currently
undergoing changes that we do not explicitly address. As detailed in Wolla
(2019), in an environment where there are ample reserves, the usual analysis
where the supply and demand for reserves determine an equilibrium federal
funds rate does not apply. The Fed now pays interest on both required and
excess reserves the sets on upper and lower bound on the federal funds rate.
From a macroeconomic modeling perspective, whether these considerations are
an innocuous bit of plumbing or factors that need explicitly modeling
remains to be seen.

Another goal of this work is to provide a guide to related outside readings
that are accessible. \ References with a (*) are suitable for
undergraduates. \ The text also includes some questions (Q:) that are
intended for pedagogical use as discussion questions or essay topics.

Section 2 discusses the TR-FR framework using graphical analysis and the
policy responses to the ZLB. Section 3 describes some associated empirical
exercises including the estimation of interest rate rules. Section 4
demonstrates algebraic analysis of a small macroeconomics model with the TR
and FR that includes solutions under rational expectations and stability
analysis under adaptive expectations. Section 5 concludes.

\section{Graphs}

The key graph combines a Taylor Rule and a Fisher Relation approximation in
the space $\left( \pi ,i\right) $ where the inflation rate is $\pi $, and
the nominal interest rate is $i$. \ As preparation, students should have
studied a macroeconomic model that describes short run fluctuations such as
the IS-LM, AS-AD-PC graphs common in most intermediate level textbooks. They
should understand that the policymaker cannot target the money supply and an
interest rate simultaneously.

A good way to introduce interest rate rules is to discuss stability under
money supply versus interest rate pegs with a money supply and demand graph,
summarizing Friedman (1968). Under an interest rate peg, a shock to money
demand would require a change in the money supply to maintain the peg that
would exacerbate the shock, shifting demand again. On a money supply and
demand graph, shifts in demand are followed by shifts in supply
indefinitely, representing an inflationary or deflationary spiral.
Alternatively, one can use Wicksell's (1898) description of the \textit{%
cumulative process} that arises when the interest rate is pegged at a value
different from the natural rate, see Woodford (Section 1.3, 2003) for a more
detailed discussion.

Q: Why don't Taylor rules suffer from the same instability?

A: The interest rate responds to changes in the economy\footnote{%
How strongly they must respond is addressed in the section on algebraic
formulations of the model.}.

Assume that the nominal rate is the policy rate targeted by the central
bank. A simple Taylor Rule takes the form%
\begin{equation}
i_{t}=\pi _{t}+\widetilde{r}+b_{\pi }\left( \pi _{t}-\bar{\pi}\right)
+b_{y}\left( y_{t}-\bar{y}\right) ,  \label{TR}
\end{equation}%
where the equilibrium real rate is $\widetilde{r}$, the inflation gap is $%
\pi _{t}-\bar{\pi},$ the output gap is $y_{t}-\bar{y}$ and the parameters $%
b_{\pi }$ and $b_{y}$ reflect the relative concern the policymaker has for
the two goals of smoothing inflation and output. To represent the Taylor
Rule on the graph, make the extra assumption that $b_{y}=0$, which means the
the monetary policymaker is engaged in pure inflation targeting, so the rule
becomes \ 
\begin{equation}
i_{t}=\pi _{t}+\widetilde{r}+b_{\pi }\left( \pi _{t}-\bar{\pi}\right) .
\label{TR infl}
\end{equation}%
Besides simplifying the pedagogy, the assumption is reasonable as it is part
of the mandate of the European Central Bank\footnote{%
The Federal Reserve has dual mandate to stabilize employment as well, which
can be accomplished by smoothing real GDP, which is represented in the
algebraic model in Section 4.}.

The Fisher Relation approximation is standard%
\begin{equation*}
r_{t}=i_{t}-E_{t}\pi _{t+1},
\end{equation*}%
where inflation expectations are $E_{t}\pi _{t+1}$. \ Making the additional
assumptions that expectations are formed adaptively such that $E_{t}\pi
_{t+1}=\pi _{t}$, following Mankiw (2013) and Buttet and Roy (2014), and
that the real rate is stable at $\widetilde{r},$\ the relation used for the
graph is%
\begin{equation*}
\widetilde{r}=i_{t}-\pi _{t}.
\end{equation*}%
A macroeconomic theory paper\footnote{%
Evans, Guse and Honkapohja (2008) and Schmitt-Grohe and Uribe (2009) are two
eamples of theory papers that include interest rate rules, the Fisher
relation and the ZLB.} might focus on steady states instead of resorting to
adaptive expectations, and some instructors may feel more comfortable with
such an approach. \ Here, for ease of exposition, it is sufficient to
describe inflation and the nominal rate as fluctuating every period, while
the equilibrium real rate $\widetilde{r}$ is determined by many factors such
as savings preferences and the perceived return on capital, which change
infrequently. The algebraic formulations of the model (Section 4) do not
make such an assumption, and the role of expectations is explored in detail.

When introducing the graph, the following questions can stimulate a brief
discussion that gives some intuition.

Q: \ What are the slopes and intercepts of the two lines?

Q: \ Where and how do they intersect?

\begin{figure}[ht]
\centering
\includegraphics[width=4.0421in]{RBQA8C00.png}
\caption{TR-FR graph}
\label{fig:trfr}
\end{figure}

In Figure 1, the vertical intercept of the FR line corresponds to the
equilibrium real rate $\widetilde{r}$, since it satisfies the Fisher
Relation when inflation is zero.\ The slope of the FR line is one, while the
slope of the TR is $1+b_{\pi }$. Assuming that the parameter $b_{\pi }$ is
positive, which corresponds to the Taylor principle, the TR is steeper than
the FR. \ Examining a time series graph of inflation and the federal funds
rate shows that the Taylor principle has usually been met, but not always,
as during the 1970s. The issue is analyzed formally in the empirical and
algebraic formulations of the model in Section 3.

The TR and FR intersect at the point $\left( \bar{\pi},\widetilde{r}+\bar{\pi%
}\right) $. \ If policy is functioning well, the data should fluctuate
around this equilibrium. \ As inflation rises above target, the policy rate
increases by a greater amount so that the real rate rises as well, and the
resulting dampening of demand should bring inflation back toward target. \
But the story does not end here.

Q: \ According to the graph, what happens when inflation is close to zero?

The problem evident from the construction of the graph in Figure 1 is that
for sufficiently low levels of inflation, the TR recommends negative values
for the policy rate. Given current institutional arrangements, negative
interest rates are not practical.\ Hence the zero lower bound on interest
rates binds, and the TR must have a kink at the horizontal axis, as in
Figure 2.\ In practice, a small negative value for short term rates might be
possible, but for simplicity, setting the bound at zero is sensible.

\begin{figure}[ht]
\centering
\includegraphics[width=4.0421in]{RBQA8C01.png}
\caption{TR-FR graph with ZLB}
\label{fig:zlb}
\end{figure}

The following gives students some intuition about the ZLB.

Q: \ What would you do if your checking account switched to pay -10\%
interest?

A: \ Naturally, everyone would withdraw their money since the 0\% interest
paid on cash is much better than -10\%. \ 

\section{ZLB on standard graphs}

Before exploring the implication of the ZLB on the TR/FR graph, it is useful
to discuss the implications for intermediate level graphs. This material is
standard, mostly. \ A number of intermediate level textbooks discuss the
liquidity trap by showing a kinked or non-linear LM curve and/or money
demand curve, for example Blanchard (2021).\ Those are included here along
with a loanable funds graph that motivates the IS relation, see Mankiw
(2013), and gives intuition about the ZLB.

Q: \ What is the economic reason for savings to equal investment $\left(
S=I\right) $? \ What would happen if that wasn't true $\left( S>I\right) $?

The loanable funds graph embodies the reason for achieving equilibrium. \ If
savings was greater than investment, bank reserves would build up, and banks
would lower interest rates in response. \ If the ZLB binds on the loanable
funds graph, as in Figure 3, it acts as a price floor, so there is excess
savings.

\begin{figure}[ht]
\centering
\includegraphics[width=4.0421in]{RBQA8C02.png}
\caption{Loanable funds graph}
\label{fig:loanable}
\end{figure}

On the corresponding IS-LM graph in Figure 4, showing a \textit{liquidity
trap}, where the IS curve has shifted down, as is typical in a recession,
the equilibrium rate is negative. However, the level of output $Y^{\ast }$\
achieved is determined by the intersection of the IS curve and the ZLB,
since negative rates are unachievable. \ The natural rate of output is shown
as $Y_{n}$.

One subtle point concerning the representation of the liquidity trap is that
the (flow of) excess savings on the loanable funds graph in Figure 3 is
temporary, since output falls and shifts savings to the left to restore
equilibrium, albeit at a level of income below the natural rate.

\begin{figure}[ht]
\centering
\includegraphics[width=4.0413in]{RBQA8C03.png}
\caption{IS-LM liquidity trap}
\label{fig:islm}
\end{figure}

The kink in the LM curve can be motivated by the kink in money demand due to
the ZLB (see Blanchard 2021).\ A number of points are evident about a
liquidity trap from the above graph.

\begin{itemize}
\item Monetary policy is ineffective, since shifts in LM do not change $%
Y^{\ast }$.

\item Fiscal policy is effective, since shifts in IS do change $Y^{\ast }$.

\item It is possible to increase $Y^{\ast }$ without crowding out with
fiscal policy, since the interest rate remains at zero.
\end{itemize}

The liquidity trap is a primary counter-example to Say's Law, which implies
that gluts, such as a sharp rise in unemployment, are impossible, since
prices can adjust to clear all markets ("supply creates its own demand"). \
As Mill (1844) pointed out and Keynes (1935) discussed at length, that logic
does not apply if some consumers refuse to spend their savings. \ A negative
equilibrium interest rate with the ZLB reflects such a state without an
obvious market solution\footnote{%
Keynes (1935, Chapter 16 section III) said that eventually capital would
depreciate simulating investment and raising the equilibrium interest rate,
but that solution could be protracted.}.

\subsection{Policy Responses}

Armed with these tools (graphs), students can better understand the policy
responses that have been employed when the ZLB is a concern, particularly in
response to the recession of 2008-9, and a couple that were proposed but not
implemented. \ Two major parts of the policy response in the U.S. can be
represented on IS-LM graph. \ Quantitative easing (QE) is an increase in the
money supply that shifts LM to the right, and the ARRA legislation in 2009
was a combination of the spending increases and tax cuts that shifts IS to
the right. In principle, the combination of the two policies could close the
recessionary gap with no crowding out from an increase in the interest rate.
\ 

Forward guidance, the promise to keep the policy rate near zero into the
future, was the other main policy response. \ It could be represented as an
increase in inflation expectations on the loanable funds and IS-LM graphs
where the nominal rate is on the vertical axis. One can argue that an
increase in expected inflation increases the demand for money and shifts LM
right, while it lowers the real rate, which shifts IS to the right due to
the impact on savings and investment decisions. \ Both shifts have the
effect of increasing output. Of course, this argument requires a more
detailed description of the IS-LM model than is common in many intermediate
levels textbooks. It also requires treating expected inflation as an
exogenous variable under the control of the monetary policymaker. Not
everyone finds such an assumption palatable, though it does point up the
importance of credibility to the effectiveness of forward guidance.

The narrative argument about the impact of forward guidance on long term
rates may be the most important point to make. A credible promise to keep a
short term rate near zero, should lower longer term rate according to the
expectations hypothesis, and those rates are the most relevant to investment
decisions. \ 

Similarly, there are potential effects of QE not represented by a shift in
LM. Alternative channels included inflation expectations, the stock market
and the foreign exchange market.

\subsubsection{Policies untried}

Two policies that were not implemented have a natural interpretation on the
TR-FR graph, the Neo-Fisherian policy of raising the lower bound on policy
rates and raising the inflation target. \ The latter shifts the TR function
to right, moving the equilibrium out along the FR. \ The appeal is that
there is more room between the equilibrium and the ZLB, which is thereby
less likely to be a factor. \ However, such a change would be difficult to
make in the middle of a crisis, and is best done as a preventative measure.
\ Krugman (1998) argues that a credible promise to increase the money supply
could accomplish it, but once the ZLB has been achieved, that credibility
could be in short supply. Whether such a change is politically feasible is
questionable as well. Banks in the business of making fixed rate home loans
would likely object, as higher inflation would lower the real return.
Theoretical work including the ZLB such as Eggertson, Mehrotra and Robbins
(2019).point out the danger of an insufficiently large increase in the
inflation target, which would generate all the costs of higher inflation
with none of the benefits.

The Neo-Fisherian solution to the ZLB is elegant in its simplicity, but
flies in the face of standard intuition about the role of money. The logic
is that, if interest rates at zero is the problem, one should raise the
bound in the interest rate rule so that the kink in the TR function would
occur at a positive value, as in Figure 5. \ Hence, the equilibrium
determined by the bound and the intersection with the FR occurs at a higher
level of inflation.

\begin{figure}[ht]
\centering
\includegraphics[width=4.0421in]{RBQA8C04.png}
\caption{Neo-Fisherian policy}
\label{fig:neofisher}
\end{figure}

The usual way the ZLB is encountered is a fall in demand leading to a
recession. To raise interest rates in such a situation is the opposite of
the usual recommendation of a Taylor Rule.

\subsubsection{Multiple equilibria and stability}

A deeper issue is the stability of the equilibria where the TR and FR
intersect. \ There are mathematically sophisticated models to address this
such as the expectational stability analysis of Evans and Honkapohja (2001,
2003,2006), see sections 4.2.2 and 4.2.3 for a simple example, but the
intuition goes back to our discussion about the instability of interest rate
pegs, motivating Taylor Rules, at the beginning of section 2. \ 

That instability argument applies to the Neo-Fisherian policy as well, where
raising the interest rate peg would require a decrease in the money supply
with a resulting deflationary spiral. The policy does not appear to have
been seriously considered by any central bank\footnote{%
See Evans and McGough (2018) for a discussion of the related theory.
Recently, Turkey tried to lower inflation by lowering interest rates with
unfortunate results. \ }.\ However, the instability issue is apparent
without any increase in the lower bound.

The equilibrium at the intersection between the FR line and the ZLB portion
of the TR function is a topic of much discussion within macroeconomics. \
The instability that could arise from a pegged interest rate is relevant at
zero. \ However, as Bullard (2010*) notes, data from the lost decade (and
beyond) for Japan seems to fluctuate around such an equilibrium in a way
that suggests stability. Figure 6 is an updated version of the first graph
in Bullard (2010*). \ Furthermore, mathematical studies with the ZLB can
show different stability results depending on the details of the model%
\footnote{%
For example, see Eusepi and Preston (2007), Benhabib, Schmitt-Grohe and
Uribe (2002), Evans, Guse and Honkapohja (2008) and Michaillat and Saez
(2021).}.

\begin{figure}[ht]
\centering
\includegraphics[width=4.0413in]{RBQA8C05.png}
\caption{U.S. and Japan data}
\label{fig:data}
\end{figure}

So there are not definitive answers about the interpretation and relevance
of the second equilibrium, but that only adds to the interest. \ Questions
for discussion include:

\begin{itemize}
\item Does the Japanese data really give evidence of a second equilibrium?

\item Is the U.S. or has the U.S. ever been in danger of approaching such a
state?

\item Does the lack of a deflationary spiral in the data mean that the
second equilibrium is meaningless?

\item Does recent data affect any of these judgements? \ 
\end{itemize}

The discussion motivates the empirical exercises discussed below.

\subsection{Secular Stagnation}

Secular stagnation has implications for monetary policy that are apparent on
the TR-FR graph. \ There are two related versions of secular stagnation,
long term declines in productivity growth and real interest rates (VoxEU
2014*). \ We can interpret the latter as a fall in the parameter $\widetilde{%
r}$ on the graph.

Students should be able to answer the following.

Q: \ If $\widetilde{r}$ falls, how does the TR-FR graph change?

A: \ Both the FR and the sloped portion of the TR shift down.

The result is the "good" equilibrium also shifts down toward the ZLB, making
the bound more likely to bind. \ A more specific version of the question
above is the following.

Q: If the value of $\widetilde{r}$ falls below $-\bar{\pi}$, what is the
problem for monetary policy?

In this case the equilibrium at the inflation target $\bar{\pi}$ requires a
negative nominal rate, so it is unachievable\footnote{%
See Appendix 7.1 and Waters (2021) for other graphing problems.}.\ Such a
fall might seem implausible, but negative real rates have been common in
developed countries for the last couple decades.

If students are familiar with the Solow model, they can make the connection
between the two flavors of secular stagnation. As productivity growth falls,
the steady state intensive capital level rises so the steady state real
interest rate falls under decreasing marginal returns.

\subsection{More Policies Untried}

Concerns about secular stagnation and the sluggish recoveries in Japan
1991-2011 (at least) and in the U.S. 2009-2014 suggest that the ZLB will be
a concern for macroeconomic policy for some time to come. \ Besides the
Neo-Fisherian policy and raising the inflation target, proposed policies
that have not yet been implemented include targeting longer term bond yields
and implementing negative rates.

Even through the ZLB period in the U.S., longer term rates such as the ten
year bond yield remained significantly positive, above 1\% with the
exception of the first few months at the beginning of the COVID pandemic
where they approached 0.5\%.\ Brunner and Meltzer (1968) argued that the
presence of positive long term rates means that liquidity traps have never
occurred. Flow of funds data from the years following the housing crisis
make this judgment hard to accept. But there is the possibility that
lowering longer term rates could stimulate the economy in such a situation.
A central bank can lower any bond yield as far as it wants by purchasing a
sufficient quantity of that bond\footnote{%
The Federal Reserve purchased long term bonds as part of QE and Operation
Twist in 2011-2 but did not explicitly target long term yields.}. In
principle, monetary policy could flatten the yield curve near zero for all
rates. \ Whether such a policy would provide sufficient stimulus to exit a
liquidity trap is uncertain, but the cost seems low. \ 

Rogoff (2017*) discusses a number of methods for removing the institutional
barriers to negative rates that would eliminate the ZLB. If the payments
system was fully electronic, negative interest rates would be quite
possible, but if people can hold cash, the ZLB exists. Some ideas such as
that of Gesell (1916) requiring cash to be stamped so that older notes would
be valued less are impractical, though interesting for students to
contemplate. \ Another proposal is to eliminate large denomination notes,
making cash hoarding more costly. Countries such as India and North Korea
have adopted such a policy, though not with an explicit desire to achieve
negative rates.

The most practical proposal discussed by Rogoff (2017*) is breaking the
one-to-one exchange rate between cash and bank reserves. For example, if a
bank sets a -3\% rate for deposits, then cash deposited a year later would
be credited \$0.97 on the dollar. \ Two years later a one dollar deposit
would be credited \$0.9409, and so on. \ Therefore, the negative interest
rate would impact cash as well as electronic currency. Hence, monetary
policymakers could set negative rates according to an interest rate rule
without a lower bound.

\section{Algebraic Analysis}

Algebraic models are helpful for both empirical and theoretical analysis.
Small macroeconomic models including Taylor Rules can demonstrate the impact
of the choice of rule on the responses to different shocks. Stability
analysis of such models formalizes and extends Friedman's (1968*) discussion
about interest rate pegs for a broad class of interest rate rules. Whether
the interest rate rules used in practice are reasonable according to these
theoretical considerations is an important empirical question.

\subsection{Empirical Exercises}

\subsubsection{Graphing the data}

Even without estimation there are graphing exercises students can do to
familiarize themselves with the data and the underlying issues. In Figure 6,
the possible relevance of the equilibrium at $\left( i,\pi \right) =\left(
-1,0\right) $ for the Japanese data is apparent, as is the "good"
equilibrium $\left( i,\pi \right) =\left( 3,2\right) $ for the U.S.
Furthermore, some U.S. data shows drift toward the ZLB. Students can create
such graphs with updated samples or data from other countries to promote
discussion as well as refining their spreadsheet skills.

\subsubsection{Taylor Rule simulations}

More formal methods can sharpen the analysis. \ Simulating the interest rate
recommended by a given Taylor Rule (\ref{TR}) can explicitly show whether
such rules represent actual policy. Figure 7 shows an example with the
parameters $b_{\pi }=b_{y}=0.5$ set to the values in the rule first proposed
in Taylor (1993*). \ The inflation target is set to $\bar{\pi}=2\%$, the
equilibrium real rate is $\widetilde{r}=2\%$ and the output gap $y_{t}-\bar{y%
}$ is determined by the percentage difference between real GDP and potential
GDP as given by the CBO. \ Naturally, there are many alternative
specifications that could form the basis for student exercises or empirical
projects.

\begin{figure}[ht]
\centering
\includegraphics[width=4.0413in]{../TRrec.png}
\caption{Federal funds rate vs Taylor rule}
\label{fig:trrec}
\end{figure}

The federal funds rate follows the Taylor Rule recommendation much of the
time, but there are some notable exceptions. Again, the late 1970s stands
out as a time that policy was much more dovish than recommended.
Interestingly, there are other such periods, which did not result in high
inflation.\ The data following the housing crisis/recession of 2008 reflects
forward guidance as the federal funds rate stays near zero long after the
rate simulated from the Taylor Rule. \ 

Of course, all of these points could be sensitive to the particular
specification of the Taylor Rule in terms of the both the parameters $b_{\pi
},b_{y}$ and $\widetilde{r}$ and the form of the output gap. \ The output
gap could be constructed with alternative forms of detrending such as the
methods of Hodrick-Prescott or Hamilton (2018), or with growth rates of real
GDP. \ Some interest rate rules substitute for the output gap using the
difference of the unemployment rate from its natural rate. The list goes on,
which provides many opportunities for students to explore the data.

\subsubsection{Estimation}

An even more formal approach is to estimate the parameters of an interest
rate rule to reveal the preferences of the policymaker. \ The work of Judd
and Rudebusch (1998*) (JR) is an invaluable model for such exercises. \ In
particular, Taylor Rules such as equation (\ref{TR}) above are usually
mis-specified for estimation, since they do not address the dynamics
involved in the conduct of monetary policy. Central bankers often engage in
interest rate smoothing, meaning they avoid drastic changes in the target
rate. JR also add a lagged output gap term to the Taylor Rule to allow for a
greater range of policymaker preferences. \ Hence their model is as follows,
where $\widetilde{\pi }_{t}$ and $\widetilde{y}_{t}$ are the inflation and
output gaps, and the target rate is $i_{t}^{T}$.%
\begin{eqnarray*}
i_{t}^{T} &=&\pi _{t}+\widetilde{r}+b_{\pi }\widetilde{\pi }_{t}+b_{y}%
\widetilde{y}_{t}+b_{L}\widetilde{y}_{t-1} \\
\Delta i_{t} &=&\gamma \left( i_{t}^{T}-i_{t}\right) +\rho \Delta i_{t-1}
\end{eqnarray*}%
The parameter $\gamma $ represents the speed at which the policy rate $i_{t}$
approaches the target $i_{t}^{T}$, and the term $\rho \Delta i_{t-1}$
reflects a momentum effect.

Hence, the estimation equation is achieved by substituting for $i_{t}^{T}$.%
\begin{equation}
\Delta i_{t}=\gamma \widetilde{r}-\gamma i_{t}+\gamma \left( 1+b_{\pi
}\right) \widetilde{\pi }_{t}+\gamma b_{y}\widetilde{y}_{t}+\gamma b_{L}%
\widetilde{y}_{t-1}+\rho \Delta i_{t-1}  \label{JR}
\end{equation}%
JR estimate the $b$'s for the different Fed chairs Burns, Volker and
Greenspan and demonstrate the differences in their approaches. \ Many
advanced undergraduates have the ability to collect the data and conduct the
estimation, which uses ordinary least squares, though there are some subtle
issues in the interpretation, about which JR is an excellent guide.
Estimating this equation is a case where entering the equation in the
statistics software, as opposed to a list of variables, is helpful.

JR's results confirm the received wisdom that Burns was insufficiently
aggressive in dealing with the inflation of the 1970s, though there could be
other factors. Nixon can get some blame, see Abrams (2006*) for a detailed
discussion that includes evidence from recorded conversations. \ More
recently, Drechsler, Savov and Schnabl (2020) argue that Regulation Q was a
much greater culprit behind inflation than is commonly supposed as the
return on deposits could not respond to changes in monetary policy. \ The
paper has a detailed empirical analysis of local deposit and inflation data.
\ The logic in that paper also extends to explain the low inflation during
the period 2009-2015 under the ZLB, which acted as a floor, as opposed to a
ceiling with Regulation Q.

One complication is that estimations like those in JR do not work with data
from the ZLB period, since there is almost no variation in the federal funds
rate.\ I've tried to use longer maturity rates as a proxy for the policy
stance with limited success.

\subsection{Theoretical analysis}

Some algebraic versions of a macroeconomic model including an interest rate
rule and Fisher relation can be solved under rational expectations or used
to examine dynamic stability under adaptive expectations. The material that
follows requires a bit more mathematical sophistication on the part of
students than would be necessary for a standard intermediate macroeconomics
class, though the mathematical tools can be taught along with the economics.

A common model in intermediate level books such as Blanchard (2021) or
Carlin and Soskice (2015) includes an IS relation and an aggregate supply
equation to create a small macroeconomic model similar to the following. \ 
\begin{eqnarray}
y_{t} &=&\bar{y}-c\left( r_{t}-\widetilde{r}\right) +u_{t}\text{ \ \ \ \ (IS)%
}  \label{smm} \\
y_{t} &=&\bar{y}+a\left( \pi _{t}-E_{t-1}\pi _{t}\right) +v_{t}\text{ \ \ \
\ (AS)}  \notag \\
i_{t} &=&\pi _{t}+\widetilde{r}+b_{\pi }\left( \pi _{t}-\bar{\pi}\right)
+b_{y}\left( y_{t}-\bar{y}\right) \text{ \ \ \ \ (TR)}  \notag \\
r_{t} &=&i_{t}-E_{t}\pi _{t+1}\text{ \ \ \ \ (FR)}  \notag
\end{eqnarray}

The IS, TR and FR equations can be solved to give an aggregate demand
relation in $\left( y,\pi \right) $ space. The expectations formulation
gives a \textit{New Classical version of} the AS equation, while supply
relations with a forward looking term $E_{t}\pi _{t+1}$ are referred to as 
\textit{New Keynesian} and are standard for much monetary policy analysis.
The form adopted here is chosen to present multiple issues in a concise way.

\subsubsection{Solutions under rational expectations}

Solutions to the model (\ref{smm}) under rational expectations demonstrate
the implications of the choice of Taylor rule parameters $b_{\pi }$ and $%
b_{y}$ representing policymaker preferences in the presence of supply and
demand shocks. We also present a solution at the ZLB, where policy in
constrained.

The method of undetermined coefficients of McCallum (1983) is the most
common and accessible solution technique. Let the demand shocks for the
model (\ref{smm}) be zero $u_{t}=0$, to focus on the case with mean zero
supply shocks $v_{t}$. To use the method of undetermined coefficients,
postulate that inflation has a solution of the form%
\begin{equation*}
\pi _{t}=d_{0}+d_{1}v_{t},
\end{equation*}%
for some coefficients $d_{0},d_{1}$, see the Appendix for details.
Substituting for the nominal rate $i_{t}$ in the IS relation using the
Fisher relation (FR) and Taylor rule (TR), yields an aggregate demand (AD)
relation.%
\begin{equation}
-c\left[ \left( 1+b_{\pi }\right) \pi _{t}-b_{\pi }\bar{\pi}-E_{t}\pi _{t+1}%
\right] =\left( 1+cb_{y}\right) \left( y_{t}-\bar{y}\right)  \label{AD RE}
\end{equation}%
Note that under the postulated solution, expected inflation is a constant $%
E_{t}\pi _{t+1}=d_{0}$, assuming the shocks $v_{t}$ are mean zero. \ Hence,
one can interpret the value%
\begin{equation}
b_{AD}=-\frac{1+cb_{y}}{c\left( 1+b_{\pi }\right) }  \label{bAD}
\end{equation}%
as the slope of AD, the avenue though which the preferences of the monetary
policymaker enter the model.

Substituting for $y_{t}-\bar{y}$ using the aggregate supply relation (AS)
and the postulated solution for inflation yields the following.%
\begin{eqnarray*}
y_{t} &=&\bar{y}+\dfrac{c\left( 1+b_{\pi }\right) }{c\left( 1+b_{\pi
}\right) +a\left( 1+cb_{y}\right) }v_{t} \\
\pi _{t} &=&\bar{\pi}-\dfrac{1+cb_{y}}{c\left( 1+b_{\pi }\right) +a\left(
1+cb_{y}\right) }v_{t}
\end{eqnarray*}%
\ The trade-off faced by policymakers dealing with a supply shock is evident
from the different signs on the shock so that output and inflation move in
opposite directions. To show the role of policymaker preferences, one can
write the solution for output using the expression for the slope $b_{AD}$
from (\ref{bAD}).%
\begin{equation*}
y_{t}=\bar{y}+\frac{1}{1-ab_{AD}}v_{t}
\end{equation*}%
Q: How would an increase in $b_{y}$ affect the response of output and
inflation to a decline in oil prices (a positive supply shock)?

As the value $b_{y}$ increases, so does the magnitude of slope of AD ($%
b_{AD}<0$ gets farther from zero) , which minimizes the impact of a supply
shock on $y_{t}$, which is apparent from fall in the the coefficient on $%
v_{t}$ in the solution for $y_{t}$ above. An increase in $b_{y}$ means the
policymaker cares more about smoothing output, and the result shows that $%
y_{t}$ is more insulated from the shock.

The solution above is the \textit{minimum state variables }solution, and
there may be others. The existence of sunspot or bubble equilibria may be
intriguing to undergraduates though formal analysis is beyond the goals of
most classes, see for example Evans and McGough (2005).

One can also solve under the assumption of a pegged nominal rate $\bar{\imath%
}$, which includes the ZLB as a special case $\bar{\imath}=0$. Here, the TR
in the model (\ref{smm}) is irrelevant and, using the FR, the IS relation
becomes%
\begin{equation*}
y_{t}=\bar{y}+cE_{t}\pi _{t+1}-c\left( \widetilde{r}+\bar{\imath}\right) .
\end{equation*}%
Using the same minimum state variable techniques as above gives the
following solutions.%
\begin{eqnarray*}
y_{t} &=&\bar{y} \\
\pi _{t} &=&-\widetilde{r}+\bar{\imath}-a^{-1}v_{t}
\end{eqnarray*}%
As in Buttet and Roy (2014), equilibrium output remains at target under the
ZLB, but inflation is typically below target. At the ZLB $\bar{\imath}=0$,
inflation fluctuates around the equilibrium represented by the intersection
of the FR and the ZLB on Figure 2.

Q: Could the solution above be used to argue for or against Neo-Fisherian
policy?

The logic of the Neo-Fisherian policy recommendation is apparent since
raising the interest rate peg $\bar{\imath}$ lifts the rational expectations
equilibrium level of inflation to any desired level. However, analysis of
stability under learning raises questions about such a policy, as noted in
Friedman (1968).

\subsubsection{Stability}

The stability analysis here formalizes the earlier discussions about
interest rate pegs and the Taylor principle. The simplest version of
learning is the adaptive expectations used earlier in the development of
Figure 2.%
\begin{equation*}
E_{t}\pi _{t+1}=\pi _{t}
\end{equation*}

Such a simple version of adaptive expectations has limitations. For example,
with New Keynesian timing of expectations in the AS equation, the inflation
terms cancel and AS is vertical. Such a result reflecting the classical
dichotomy would not arise with almost any other version of learning.
Expectational stability of Evans and Honkapohja (2001) is the formulation of
learning that is most generally applicable for stability analysis. The form $%
E_{t}\pi _{t+1}=\pi _{t}$ is chosen here strictly for ease of exposition.

Using this assumption in the AS relation in the model (\ref{smm}), and
aggregate demand relation achieved by substituting for the variables $%
r_{t},i_{t}$ with the other three equations in the model (\ref{smm}) results
in the following.%
\begin{eqnarray*}
y_{t} &=&\bar{y}-c\left( \bar{\imath}-\pi _{t}-\widetilde{r}\right) \text{ \
\ \ \ (AD)} \\
y_{t} &=&\bar{y}+a\left( \pi _{t}-\pi _{t-1}\right) \text{ \ \ \ \ (AS)}
\end{eqnarray*}

Buttet and Roy (2014) interpret the model such that the slopes of AD and AS
are $c$ and $a$ respectively. The fact that AD is upward sloping is an
indication that the interest rate peg has perverse effects. The
interpretation of the slopes is intuitive but different than under rational
expectations where expected inflation is a constant. \ Here, expected
inflation becomes the endogenous level of inflation in the AD relation and
an exogenous lagged value in AS.

Solving to eliminate output from the AD and AS equations above yields the
following difference equation.%
\begin{equation}
\pi _{t}=\frac{a}{a-c}\pi _{t-1}+\frac{c}{a-c}\left( \widetilde{r}-\overline{%
i}\right)  \label{diff peg}
\end{equation}

Reasonable values of the parameters show instability, meaning inflation
diverges. \ Both parameters $a$ and $c$ are positive and inflation shows
persistence in the data so the fraction $\dfrac{a}{a-c}$ should also be
positive, which implies $a>c$. \ In that case, it is indeed the case that $%
\dfrac{a}{a-c}>1$, and the equilibrium is unstable under adaptive
expectations for a fixed interest rate policy, which corresponds to the
discussion of a pegged rate around Figure 5. \ 

There is one solution to the difference equation where the initial value $%
\pi _{0}$ is at the steady state value $\pi ^{\ast }=\widetilde{r}-\bar{%
\imath}$. \ However, policymakers do not get to choose their initial values,
and the inclusion of any shock would push inflation onto a divergent path. \
Whether it diverges to positive or negative infinity depends on the initial
value $\pi _{0}$ and the peg $\bar{\imath}$. If the initial value of
inflation is less than the steady state such that $\pi _{0}<\widetilde{r}-%
\bar{\imath},$ inflation falls indefinitely, a \textit{deflationary spiral}.

\subsubsection{Taylor Principle}

For a model with a Taylor Rule, one can show the relevance of the Taylor
principle via stability analysis. Again assuming adaptive expectations $%
E_{t}\pi _{t+1}=\pi _{t}$, the model (\ref{smm}) above reduces to the
following difference equation.%
\begin{equation}
\pi _{t}=\frac{a\left( 1+cb_{y}\right) }{a\left( 1+cb_{y}\right) +cb_{\pi }}%
\pi _{t-1}+\frac{cb_{\pi }}{a\left( 1+cb_{y}\right) +cb_{\pi }}\bar{\pi}
\label{diff Taylor}
\end{equation}%
Q: \ Is there a simple condition on the policymaker preference parameters $%
b_{\pi },b_{y}$ that guarantess convergence/stability?

If the Taylor principle $b_{\pi }>0$ holds, the coefficient on $\pi _{t-1}$
is between 0 and 1 and inflation converges but, if not, inflation diverges,
demonstrating instability, a result found in Mankiw (2013). The cases where
the Taylor principle holds are examples of interest rate rules that are not
subject to Friedman's (1968*) objection concerning interest rate pegs, since
the policymaker is sufficiently aggressive in adjusting the rate in response
to economic conditions. \ Evans and Honkapohja (2003, 2006) survey results
about expectational stability and determinacy for a range of interest rate
rules. \ A related discussion can also be found in Woodford (2003, Chapter
4, section 2.3).

\subsubsection{Stability under the ZLB}

A stability analysis that encompasses the ZLB requires an interest rate rule
with a kink as in Figure 2. One example uses the following interest rate rule%
\begin{equation*}
i_{t}=\max \left\{ 0,\pi _{t}+\widetilde{r}+b_{\pi }\left( \pi _{t}-\bar{\pi}%
\right) +b_{y}\left( y_{t}-\bar{y}\right) \right\} .
\end{equation*}%
Using this interest rule in place of the TR in the model (\ref{smm}), there
are two steady states corresponding to the intersections of TR and FR on the
graph. \ One is the steady state for the difference equation (\ref{diff
Taylor}) when the ZLB does not bind, which is stable at the inflation target 
$\pi ^{\ast }=\bar{\pi}$. The other is the steady state for the difference
equation (\ref{diff peg}) for $\bar{\imath}=0$, which is unstable at the
value $\pi ^{\ast }=-\widetilde{r}$. \ Output is at its target value for
both steady states\footnote{%
A more sophisticate model such as the one in Jung, Teranishi and Watanabe
(2005) is required for there to be a distinction in the values of output at
the steady states.}.

The kinked interest rate rule above determines the condition for whether the
ZLB binds. \ The nominal rate is positive if%
\begin{equation*}
\pi _{t}+\widetilde{r}+b_{\pi }\left( \pi _{t}-\bar{\pi}\right) +b_{y}\left(
y_{t}-\bar{y}\right) >0.
\end{equation*}%
This condition is satisfied at the stable steady state where both output and
inflation are at target as long as $\bar{\pi}+\widetilde{r}>0$. \ However,
this condition shows that a sufficiently large adverse shock to the
equilibrium natural rate $\widetilde{r}$ could make the stable steady
impossible to achieve so that the the nominal rate must be zero, creating
possible instability. \ 

On the graph in Figure 8, such a fall in $\widetilde{r}$, as with secular
stagnation, is represented by a shift down in both the TR and FR relations.

\begin{figure}[ht]
\centering
\includegraphics[width=4.0421in]{RBQA8C06.png}
\caption{Secular stagnation}
\label{fig:secular}
\end{figure}

A simple way to show a fiscal policy response to such a decline in the
equilibrium real rate is to included a government spending term as a shock
to the IS equation, like the variable $u$ in the model (\ref{smm}). \ A
positive shock, meaning an increase in spending, would counter the fall in $%
\widetilde{r}$.

Buttet and Roy (2014) analyze the model above with the ZLB under adaptive
expectations and graphical stability analysis. The AD relation has a kink
such that the slope of AD is positive at $c>0$ when the ZLB binds but is
downward sloping otherwise. They also show that an equilibrium where AS
intersects the upward sloping portion of AD is unstable when the slope of AD
is greater than the slope of AS using a dynamic analysis where changes in
lagged inflation shift AD. Fiscal stimulus shifts AD to the right and can
move the equilibrium off the upward sloping portion so inflation converges
to target $\bar{\pi}$.

Similarly, given the interest rate peg $i_{t}=\bar{\imath}$, the combination
of the FR and IS relations in the model (\ref{smm}) is an upward sloping
aggregate demand relation with slope $c$. \ The condition $a>c$ that
guarantees instability for any interest rate peg (\ref{diff peg}) is
equivalent to the slope of aggregate demand being greater than the slope of
aggregate supply.

Related graphs in more sophisticated models can also be found in Eggertson
and Egiev (2019) and Eggertson, Mehrotra and Robbins (2019). Advanced
undergraduates should be able to understand the graphical stability analysis
in Buttet and Roy (2014) or the algebraic analysis presented here or both
depending on the time constraints and their mathematical background. A
version using a Taylor Rule that only responds to inflation (\ref{TR infl})
where $b_{y}=0$\ simplifies the exposition.

\subsubsection{Forward guidance}

As discussed in Section 3, an important policy response to the ZLB is
forward guidance, the promise to keep a targeted interest rate near zero for
a specified length of time in the future. The intuition behind such a policy
is that a credible promise to engage in monetary stimulus beyond normal
policy, and the resulting increase inflation expectations can stimulate the
economy today.

Unfortunately, the proper way to model forward guidance is far from settled
law in macroeconomics. The work most closely related\footnote{%
Some alternative approaches include Krugman (1998), who presents a two
period model where the central bank promises to "irresponsibly" increase
inflation in the future, though the policymaker is assumed to have that
ability. In the overlapping generations model of Eggertson, Mehrotra and
Robbins (2019), forward guidance is not effective at lifting the economy
away from the ZLB. Eggertson and Egiev (2019) describe a New Keynesian model
based on Eggertson and Woodford (2003) where there is an exogenous regime
shift in the equilibrium real rate. \ They describe the effect of
expectations as a change in the probability of the shift.} to the model (\ref%
{smm}) that includes forward guidance is Jung, Teranishi and Watanabe
(2005).\ They study the optimal response of monetary policy to a shift down
in the equilibrium real rate, as in Figure 8, followed by a gradual return
to a value where the optimal policy rate is above the ZLB. \ Forward
guidance emerges by comparing the policy response under discretion, when the
policymaker optimizes in a single period, and commitment, where the
policymaker considers the long run path of policy, the \textit{timeless
perspective} in the terminology of Woodford (2003). Under commitment, the
policymaker keeps the nominal rate at zero beyond the time that the
equilibrium real rate has returned to its original value, so forward
guidance emerges endogenously.

Discussing forward guidance with the algebraic models presented here is not
practical for most undergraduate classes. Commitment versus discretion for
monetary policymakers is interesting, but developing mathematical models
that express the timeless perspective is a topic for a Ph.D. level class.

\subsection{Existing Textbook Treatments}

Almost all undergraduate macroeconomics textbook contain elements in the
discussion above, but none of them explicitly combine the Taylor Rule and
Fisher relation graphically or algebraically. \ All of those referenced here
contain an explanation of real vs. nominal interest rates, which includes
the Fisher relation, and an algebraic expression of a Taylor Rule, with the
exception of Chugh (2015), which is focused on micro-founded DSGE models. \
Blanchard (2021), Jones (2021) and Carlin and Soskice (2015) discuss the
liquidity trap and deflation spirals by specifying a FR under the ZLB such
that $r_{t}=-\pi _{t}$ combined with a dynamic Phillips Curve. \ Blanchard
(2021) also discusses the liquidity trap/ZLB in the context of a money
supply and demand graph\footnote{%
A referee notes that forward guidance is discussed with a standard IS-LM
model in Mishkin (2015).}. Mankiw (2021) has a model closely related to the
one described here (\ref{smm}) under adaptive expectations where the TR and
IS relations are expressed in terms of the real rate, so the FR is
superfluous. \ Graphical analysis of the output-inflation trade-off and the
Taylor principle reflect the algebraic analysis here.

\section{Summary}

The trends behind secular stagnation indicate that the ZLB will be a fact of
life for decades to come, which has implications for stabilization policy in
general and the conduct of monetary policy in particular. Understanding the
motivation for and the implementation of the extraordinary policies in the
presence of the ZLB is important for many outside the economics profession.
The graphs with the TR and FR give a simple framework for organizing many of
the issues and the related data. For those who want to understand current
research on these issues, the study of solution techniques and stability
analysis in models including these two relations is an essential stepping
stone.

While addressing these policy issues, undergraduates can get experience with
empirical techniques from creating spreadsheet graphs to estimating Taylor
Rules. They can also learn to solve models under rational and adaptive
expectations. Finally, the combination of all these methodologies helps to
form connections between the data, the graphs and the algebraic models, a
goal of most economics education.\bigskip

\section{References}

\newline
Abram, Burton A. (2006) \ How Richard Nixon Pressured Arthur Burns: \
Evidence from the Nixon Tapes. \ \textit{Journal of Economic Perspectives }%
20(4), 177-188.\bigskip\ \newline
Benhabib, Jess, Schmitt-Groh\'{e}, Stephanie and Mart\'{\i}n Uribe (2002).
Avoiding Liquidity Traps. \textit{Journal of Political Economy 110(3), }p.
535-563\bigskip\ \newline
Blanchard, Olivier (2021). \textit{Macroeconomics, 8th edition.}
Pearson\bigskip\ \newline
Brunner, Karl. AND Allan H. Meltzer (1968): Liquidity Traps for Money, Bank
Credit, and Interest Rates. \textit{Journal of Political Economy, 76 (1)},
1--37.\bigskip\ \newline
Bullard, James (2010). \ Seven Faces of "The Peril." \textit{\ Federal
Reserve Bank of St. Louis Review} 92(5), 339-352.\bigskip\ \newline
Buttet, Sebastien and Udayan Roy (2014). A Simple Treatment of the Liquidity
Trap for Intermediate Macroeconomics Courses, \textit{The Journal of
Economic Education, 45(1)}, 36-55.\bigskip\ \newline
Carlin, Wendy and David Soskice, (2015). \textit{Macroeconomics: \
Institutions, Instability and the Financial Crisis.} Oxford University
Press.\bigskip\ \newline
Chugh, Sanjay (2015). \textit{Modern Macroeconomics.} MIT Press, Cambridge,
MA.\bigskip\ \newline
Drechsler, Itamar, Alexei Savov and Philipp Schnabl (2020). The Financial
Origins of the Rise and Fall of American Inflation. working paper, NYU Stern
School of Business.\bigskip\ \newline
Evans, George, and Seppo Honkapohja (2001). \ \textit{Learning and
Expectations in Macroeconomcs}. \ Princeton University Press, Princeton,
NJ.\bigskip\ \newline
Evans, George and Seppo Honkapohja (2003) Expectations and the Stability
Problem for Optimal Monetary Policies. \textit{Review of Economic Studies}
70, 807-824.\bigskip \newline
Evans, George and Seppo Honkapohja (2006) Monetary Policy, Expectations and
Commitment. \textit{Scandinavian Journal of Economics} 108, 15-38.\bigskip 
\newline
Evans, George, Guse, Eran and Seppo Honkapohja (2008). Liquidity Traps,
Learning, and Stagnation. \textit{European Economic Review, 52(8)}, p.
1438-463.\bigskip\ \newline
Evans, George and Bruce McGough (2005) Equilibrium Selection, Observability
and Backward-stable Solutions. \textit{Journal of Monetary Economics} 91,
1-10.\bigskip\ \newline
Evans, George and Bruce McGough (2018) Monetary Policy, Indeterminacy and
Learning. \textit{Journal of Economic Dynamics and Control} 29,
1809-1840.\bigskip\ \newline
Eggertson, Gauti, Neil Mehrotra and Jacob Robbins (2019) A Model of Secular
Stagnation: Theory and Quantitative Evaluation. \textit{American Economic
Journal: Macroeconomics} 11(1), 1-48.\bigskip\ \newline
Eusepi, Stephano and Bruce Preston (2007) Learnability and Monetary Policy:
A Global Perspective. \textit{Journal of Monetary Policy} 54(4),
1115-1131.\bigskip\ \newline
Friedman, Milton (1968) The Role of Monetary Policy. Presidentional address
delivered at the 80th Annual Meeting of the American Economic Association. 
\textit{American Economic Review }56(1), 1-17.\bigskip\ \newline
Gesell, Silvio (1958). \textit{The Natural Economic Order.} Peter Owen,
London.\bigskip\ \newline
Hamilton, James (2018). \ Why You Should Never Use the Hodrick-Prescott
Filter. \ \textit{Review of Economics and Statistics 100(5)},
831-843.\bigskip\ \newline
Howitt, Peter (1992). \ Interest Rate Control and Non-Convergenace to
Rational Expectations. \ \textit{Journal of Political Economy} 100(4),
776-800.\bigskip\ \newline
Jones, Charles (2021). \ \textit{Macroeconomics.} W. W. Norton \& Company,
New York, NY.\bigskip\ \newline
Judd, Kenneth and Glenn D. Rudebusch (1997). \ Taylor's rule and the Fed. 
\textit{FRBSF Economic Letter} 1998-4, 3-16.\bigskip\ \newline
Jung, Taehun, Yuki Teranishi, and Tsutomu Watanabe (2005). Optimal Monetary
Policy at the Zero-Interest-Rate Bound. \textit{Journal of Money, Credit and
Banking} 37(5), 813--35.\bigskip\ \newline
Keynes, John Maynard (1936). \ \textit{The General Theory of Employment,
Interest and Money.} Palgrave Macmillan, London.\bigskip\ \newline
Mill, John Stuart (1844). \ \textit{Essays on Some Unsettled Questions in
Political Economy.} John W. Parker, London.\bigskip\ \newline
Mankiw, Gregory N. (2013).\textit{\ Macroeconomics. 8th ed}, Worth
Publishers, New York.\bigskip\ \newline
Mankiw, Gregory N. (2021).\textit{\ Macroeconomics. 11th ed}, Worth
Publishers, New York.CMichaillat, Pascal and Emmanuel Saez (2021). Resolving
New Keynesian Anomalies with Wealth in the Utility Function. \ Review of
Economics and Statistics 103(2), 197-205.\bigskip\ \newline
McCallum, Bennett T. (1983). \ On Non-Uniqueness Rational Expectations
Models. \textit{Journal of Monetary Economics }11,139-168.\bigskip\ \newline
Mishkin, Frederic (2015). \textit{Macroeconomics: \ Policy and Practice.}
Pearson.\bigskip\ \newline
Rogoff, Kenneth (2017). \ Dealing with Monetary Paralysis at the Zero Bound.
\ \textit{Journal of Economic Perspectives 31(3), }47-66.\bigskip\ \newline
Schmitt-Groh\'{e}, Stephanie and Martin Uribe (2009). Liquidity Traps with
Global Taylor Rules. \textit{International Journal of Economic Theory, 5(1)}%
, pp. 85-106.\bigskip\ \newline
Taylor, John B., (1993). \ Discretion Versus Policy Rules in Practice. \ 
\textit{Carnegie-Rochester Conference Series on Public Policy}, vol. 39,
December,195-214.\bigskip\ \newline
Vox EU (2014). \ Secular Stagnation: Facts, Causes and Cures. \ e-book,
edited by Coen Teulings and Richard Baldwin, \
https://voxeu.org/content/secular-stagnation-facts-causes-and-cures \ 
\textit{accessed 10/4/2019}.\bigskip\ \newline
Waters, George (2021). \ \textit{A Primer on Macoeconomics}.
manuscript.\bigskip\ \ \newline
Wolla, Scott (2019). \ A New Frontier: \ Monetary Policy with Ample
Reserves. \textit{Page One Economics\registered }, May 2019, The Federal
Reserve Bank of St. Louis.\bigskip\ \newline
Woodford, Michael (2003). \ \textit{Interest and Prices}. Princeton
University Press, Princeton, NJ.

\section{Appendix}

\subsection{\protect\Large Sample TR-FR Graphing Problems}

\begin{enumerate}
\item Draw the graph with the Fisher Relation and the Taylor Rule assuming
that the steady state real rate is lower than $-\bar{\pi}$, where the
inflation target is $\bar{\pi}$. \ Why would this be a particularly
difficult situation for a central bank? \ What could the Fed do to alleviate
the problem?

\item Starting from the TR-FR graph in Figure 2, show the effect of an
increase in the inflation target $\bar{\pi}$. \ How are the equilibria
affected?
\end{enumerate}

See Waters (2021) for other examples.

\subsection{Model Solution under Rational Expectations}

The model with supply shocks $v_{t}$%
\begin{eqnarray}
y_{t} &=&\bar{y}-c\left( r_{t}-\widetilde{r}\text{ }\right) \text{\ \ \ \
(IS)} \\
y_{t} &=&\bar{y}+a\left( \pi _{t}-E_{t-1}\pi _{t}\right) +v_{t}\text{ \ \ \
\ (AS)}  \notag \\
i_{t} &=&\pi _{t}+\widetilde{r}+b_{\pi }\left( \pi _{t}-\bar{\pi}\right)
+b_{y}\left( y_{t}-\bar{y}\right) \text{ \ \ \ \ (TR)}  \notag \\
r_{t} &=&i_{t}-E_{t}\pi _{t+1}\text{ \ \ \ \ (FR)}  \notag
\end{eqnarray}%
can be solved under rational expectations with the method of undetermined
coefficients. \ Substituting out $r_{t}$ using the IS, FR and TR equations 
\begin{equation*}
y_{t}-\bar{y}=.-c\left[ \pi _{t}-E_{t}\pi _{t+1}+b_{\pi }\left( \pi _{t}-%
\bar{\pi}\right) +b_{y}\left( y_{t}-\bar{y}\right) \right] ,
\end{equation*}%
gives the aggregate demand relationship%
\begin{equation}
y_{t}-\bar{y}=\frac{-c}{1+cb_{y}}\left[ \pi _{t}-E_{t}\pi _{t+1}+b_{\pi
}\left( \pi _{t}-\bar{\pi}\right) \right]  \label{ADapp}
\end{equation}%
Combining this AD relation with the AS relation substituting out $y_{t}-\bar{%
y}$ yields%
\begin{equation*}
\frac{-c}{1+cb_{y}}\left[ \pi _{t}-E_{t}\pi _{t+1}+b_{\pi }\left( \pi _{t}-%
\bar{\pi}\right) \right] =a\left( \pi _{t}-E_{t-1}\pi _{t}\right) +v_{t}.
\end{equation*}%
Postulating solutions of the forms%
\begin{equation*}
\pi _{t}=d_{0}+d_{1}v_{t}
\end{equation*}%
for unknown $d_{0},d_{1}$ and simplifying, the above relations become%
\begin{equation*}
\frac{-c}{1+cb_{y}}\left[ d_{1}v_{t}+b_{\pi }\left( d_{0}+d_{1}v_{t}-\bar{\pi%
}\right) \right] =a\left( d_{1}v_{t}\right) +v_{t}
\end{equation*}%
using the fact that $E_{t}\pi _{t+1}=E_{t-1}\pi _{t}=d_{0}$, since the shock 
$v_{t}$ is mean zero. \ 

Matching the constant terms and shock terms gives the following.%
\begin{eqnarray*}
\frac{-c}{1+cb_{y}}\left( d_{0}-\bar{\pi}\right) &=&0 \\
\frac{-c}{1+cb_{y}}\left( d_{1}+b_{\pi }d_{1}\right) &=&ad_{1}+1.
\end{eqnarray*}%
Hence, the coefficients are $d_{0}=\bar{\pi}$ and $d_{1}=-\dfrac{1+cb_{y}}{%
c\left( 1+b_{\pi }\right) +a\left( 1+cb_{y}\right) }$, so the solutions are
as in the text, using the AD relation (\ref{ADapp}) to solve for $y_{t}$.

\begin{eqnarray*}
y_{t} &=&\bar{y}+\dfrac{c\left( 1+b_{\pi }\right) }{c\left( 1+b_{\pi
}\right) +a\left( 1+cb_{y}\right) }v_{t} \\
\pi _{t} &=&\bar{\pi}-\dfrac{1+cb_{y}}{c\left( 1+b_{\pi }\right) +a\left(
1+cb_{y}\right) }v_{t}
\end{eqnarray*}%
\end{doublespace}

\end{document}
